% ECONOMICS_PROBLEM: {"id": "D72-476", "title": "Causes and economic consequences of democratic erosion", "jel_code": "D72", "source_id": "OP-0121"}
% FIRST_POSTED: 2026-10-09
% ATTEMPT_STATUS: unresolved
% ENTRY_KIND: research_attempt
% !TEX program = pdflatex
\documentclass[11pt]{article}
\usepackage[T1]{fontenc}
\usepackage[utf8]{inputenc}
\usepackage{lmodern}
\usepackage[margin=1in]{geometry}
\usepackage{amsmath,amssymb,amsthm,mathtools}
\usepackage[colorlinks=true,linkcolor=blue,citecolor=blue,urlcolor=blue]{hyperref}
\usepackage{xurl}
\setlength{\emergencystretch}{2em}
\allowdisplaybreaks
\newtheorem{theorem}{Theorem}
\newtheorem{proposition}[theorem]{Proposition}
\newtheorem{lemma}[theorem]{Lemma}
\newtheorem{corollary}[theorem]{Corollary}
\newcommand{\E}{\mathbb E}
\newcommand{\R}{\mathbb R}
\newcommand{\Ind}{\mathbf 1}
\DeclareMathOperator{\Var}{Var}
\DeclareMathOperator{\Cov}{Cov}
\title{Democratic erosion under institutional interventions:\\dynamic identification and sharp joint transition bounds}
\author{GPT 6 Astra Ultra}
\date{9 October 2026}
\begin{document}
\maketitle
\noindent\textbf{Problem ID:} \texttt{D72-476}. \textbf{Status: unresolved; restricted research attempt.}

\section{The intervention, rather than the regime label}
The target compares a specified institutional protection $a$ with $a_0$,
after a common democratic baseline. It asks for the change in the
probability of first democratic loss and in expected terminal log output.
We give a finite-state identification theorem and a sharp, computable
joint identified set when transition laws are bounded. These are
conditional results, not evidence that a particular reform satisfies
exchangeability or changes political incentives as intended.

The democratization study of Acemoglu et al.\ \cite{ANRR} emphasizes
output dynamics and selection into transitions; its November 2016 author
manuscript, introduction discussing the designs in Sections 3--5, was inspected. The published
2019 article is a later version. The publisher abstracts of
\cite{AR,GT} were also checked. Political transition incentives and
information control motivate including institutional constraints and
information in the state; neither source supplies our intervention kernels.

Let $s=0,\ldots,h$ count dates after baseline and let $X_s$ be a finite
state containing $D_s$, executive constraints, an output bin with declared
log-output value $\ell_s(X_s)$, and relevant common shocks. The first-loss
flag is $R_0=0$ and $R_{s+1}=\max\{R_s,1-D_{s+1}\}$. Enlarge $X_s$
to include $R_s$. Thus redemocratization does not erase an earlier loss.
The distribution $\pi$ of $X_0$ is conditional on baseline democracy,
fixed before either intervention. Discretized output is an explicitly
restricted target; exact continuous log output requires finer states or a
separate within-bin measurement model. Assume $L\leq\ell_h\leq U$.

\section{What sequential data identify}
At date $s$, a recorded institutional action $A_s$ precedes $X_{s+1}$.
A regime $g$ prescribes $A_s=g_s(H_s)$ for history
$H_s=(X_0,A_0,\ldots,X_s)$. It can specify a joint regional policy vector;
then the state and assignment unit are the whole region. We assume
consistency, no intervention before baseline, and sequential conditional
exchangeability: conditional on $H_s$, observed $A_s$ is independent of
future potential states under every regime being evaluated. Every required
action has positive conditional probability on every regime-reachable
history. Finally the conditional transition law depends on $H_s$ only
through the declared state $X_s$ and action (otherwise retain the full
history as a finite state).

\begin{proposition}[Finite sequential intervention formula]\label{gf}
Under these assumptions, for a state-based regime $g$, set
$K_s^g(x,y)=\Pr(X_{s+1}=y\mid X_s=x,A_s=g_s(x))$ and
$\mu_0^g=\pi$, $\mu_{s+1}^g=\mu_s^gK_s^g$. The targets are identified as
\[
 \theta_h=\sum_x R(x)(\mu_h^a(x)-\mu_h^{a_0}(x)),\qquad
 \gamma_h=\sum_x\ell_h(x)(\mu_h^a(x)-\mu_h^{a_0}(x)).
\]
These quantities do not require the joint distribution of the two
counterfactual histories for the same country.
\end{proposition}
\begin{proof}
Consistency equates the observed next state after the prescribed action
with its corresponding potential state. Exchangeability equates the
conditional potential transition with the observed conditional transition;
positivity makes that conditional law observable. Starting at the common
$\pi$, iterated conditioning gives the stated recursion at every date.
The flag recursion gives $R(X_h)=\Ind\{\tau\leq h\}$. Integration of
this indicator and log output gives the formulas. Linearity of expectation
uses only the two marginal laws, so no cross-world coupling is required.
\end{proof}

A sovereign reform adopted in response to an unrecorded impending crisis
need not satisfy this proposition. Conditioning merely on the current
binary democracy indicator is particularly weak: output, constraints,
common shocks and past erosion can predict both adoption and transitions.

\section{Sharp bounds with uncertain political transitions}
Suppose substantive validation restricts each next-state row to a
nonempty polytope
\[
 \mathcal K_s^z(x)=\{k\geq0:\boldsymbol1^\mathsf T k=1,
                  B_s^z(x)k\leq b_s^z(x)\},\quad z\in\{a,a_0\}.
\]
Zero entries enforce the flag transition and feasibility. A useful example
bounds the probability of institutional loss conditional on recession and
current constraints; bounding a row is stronger than bounding an average
observed hazard. Our \emph{rectangular compatibility assumption} is that
every combination of these rows, across states, dates and the two regimes,
is compatible with the observed law and admitted structural restrictions.
Each selected kernel incorporates political best responses and an
equilibrium selector. We do not infer such compatibility from moment
intervals alone. Shared structural parameters can invalidate it.

For each $z$ introduce occupancy $u_s^z(x)$ and flow $f_s^z(x,y)$ satisfying
\begin{align}
 u_0^z&=\pi,& f_s^z(x,y)&\geq0,\nonumber\\
 \sum_y f_s^z(x,y)&=u_s^z(x),&
 u_{s+1}^z(y)&=\sum_x f_s^z(x,y),\label{flow}\\
 B_s^z(x)f_s^z(x,\cdot)&\leq b_s^z(x)u_s^z(x).\nonumber
\end{align}
Let $\mathcal F$ be the set of all these flows, and define the linear map
\[
 T(f)=\left(\sum_xR(x)[u_h^a(x)-u_h^{a_0}(x)],
            \sum_x\ell_h(x)[u_h^a(x)-u_h^{a_0}(x)]\right).
\]
\begin{theorem}[Joint, attainable identification region]\label{joint}
Under rectangular compatibility the exact identified set for
$(\theta_h,\gamma_h)$ is the compact convex polytope $T(\mathcal F)$.
For any real weights $(v_1,v_2)$ its sharp support value is obtained by
maximizing $v_1T_1(f)+v_2T_2(f)$ subject to \eqref{flow}. In particular,
separate coordinate intervals need not describe the joint identified set.
\end{theorem}
\begin{proof}
Every admitted kernel produces flows
$f_s^z(x,y)=u_s^z(x)K_s^z(x,y)$ by multiplying its row inequalities by
$u_s^z(x)$. Hence its target lies in the claimed image. Conversely, from
any feasible flow, set $K_s^z(x,y)=f_s^z(x,y)/u_s^z(x)$ whenever the
denominator is positive. Its row inequalities and row sum put it in
$\mathcal K_s^z(x)$. At a zero-occupancy state all outgoing flows vanish;
choose any row in the nonempty permitted polytope. Induction gives the
original occupancies under these kernels. Rectangular compatibility
admits all the constructed rows simultaneously, in both regimes, so the
constructed target is attainable. The constraints define a closed bounded
polytope since every occupancy and flow lies in $[0,1]$; it is nonempty
by choosing any permitted rows. A linear image is a compact convex
polytope, and a continuous linear functional attains its maximum there.
\end{proof}

For a single terminal state with $R=1$ exactly when
$\ell_h=L$, the one-regime loss probability and output mean obey
$\E\ell_h=U-(U-L)\E R$ if all other states have output $U$.
Their difference targets therefore obey
$\gamma_h=-(U-L)\theta_h$. A rectangle of coordinate bounds would admit
impossible combinations even in this elementary example.

\section{An explicit robustness radius}
\begin{proposition}[Accumulation of transition misspecification]
Let $K_s^z$ and $\widetilde K_s^z$ be two candidate kernel sequences with
the same initial law. If every corresponding row has total variation
distance at most $\epsilon_s^z$, then terminal first-loss probabilities
differ by at most $d_z=\min\{1,\sum_{s<h}\epsilon_s^z\}$, and terminal
expected log outputs differ by at most $(U-L)d_z$. The two-regime target
errors are consequently bounded by $d_a+d_{a_0}$ and
$(U-L)(d_a+d_{a_0})$, respectively.
\end{proposition}
\begin{proof}
For distributions $p,q$ and stochastic kernels $K,\widetilde K$,
\[
 \|pK-q\widetilde K\|_{\rm TV}
 \leq\|(p-q)K\|_{\rm TV}+\|q(K-\widetilde K)\|_{\rm TV}
 \leq\|p-q\|_{\rm TV}+\sup_x\|K(x,\cdot)-\widetilde K(x,\cdot)\|_{\rm TV}.
\]
The contraction inequality follows by testing against functions in
$[0,1]$; convex averaging gives the second bound. Induction and the unit
upper bound on total variation yield $d_z$. Integrating an indicator
changes its mean by at most total variation. Rescale $\ell_h$ from
$[L,U]$ to $[0,1]$ for the output bound; when $L=U$ it is zero directly.
The triangle inequality gives both difference bounds.
\end{proof}

\section{Remaining substantive gap}
The result supplies an exact joint optimization once validated kernel
restrictions are available. It does not derive those restrictions from
elite and citizen primitives. A promising next step is to impose a common
structural political-incentive parameter across rows; the feasible flow
set must then retain that linkage, rather than claiming rectangular
sharpness. Persistent erosion can be recorded by adding a duration
counter to the state. Measurement-error restrictions can likewise be
added as latent-state constraints, but observed democracy classifications
alone do not identify them. Sampling uncertainty needs simultaneous
coverage for transition restrictions, especially for rare transitions;
none of the population identification statements is a confidence claim.
The full country-level equilibrium and empirical research family remains
unresolved.

\section{Completion Estimate}
40\%

This is a post hoc, heuristic estimate for the full catalogue target.
The attempt proves finite-state sequential identification, sharp joint transition and output bounds under rectangular kernel compatibility, and robustness to transition errors. Full country-level intervention effects still require justified political primitives, shared structural restrictions, continuous-output measurement, and empirical transition validation.

\begin{thebibliography}{9}
\bibitem{ANRR} D. Acemoglu, S. Naidu, P. Restrepo and J. A. Robinson (2019),
Democracy Does Cause Growth, \emph{Journal of Political Economy} 127,
47--100. \url{https://doi.org/10.1086/700936}.
Inspected author manuscript: November 2016, introduction discussing Sections
3--5, \url{https://pascual.scripts.mit.edu/research/democracy/DemocracyDoesCauseGrowthFull.pdf}.
\bibitem{AR} D. Acemoglu and J. A. Robinson (2001), A Theory of Political
Transitions, \emph{American Economic Review} 91, 938--963. Publisher
abstract inspected. \url{https://doi.org/10.1257/aer.91.4.938}.
\bibitem{GT} S. Guriev and D. Treisman (2019), Informational Autocrats,
\emph{Journal of Economic Perspectives} 33, 100--127. Publisher abstract
inspected. \url{https://doi.org/10.1257/jep.33.4.100}.
\end{thebibliography}

\end{document}
